\documentclass[12pt]{article}
\usepackage[utf8]{inputenc}
\usepackage[T1]{fontenc}
\usepackage{hyperref}
\title{Faifi Language Resources for K-12 Curriculum Development}
\author{Compiled by Claude Sonnet 5.5 (claude-sonnet-5-5) for Andr\'as Kornai}
\date{}
\begin{document}
\maketitle

\section{Introduction}
Faifi (Fayfi, Khawlaniyah; ISO 639-3: fif) is spoken in the Jabal Fayfa mountains of Jazan province, southwestern Saudi Arabia, by roughly 50,000 people according to Wikipedia and Joshua Project. Its classification is disputed: researchers have not settled whether it is an unusual Arabic variety or a survivor of Ancient South Arabian. It is considered endangered as younger speakers shift toward mainstream Saudi Arabic and Modern Standard Arabic.

\section{Primary Educational Resources}
No school materials in Faifi were found.

\section{Linguistic and Reference Materials}
\begin{itemize}
\item Alfaifi et al. 2025, describing Upper and Lower Faifi and giving a first acoustic description of Upper Faifi: \url{https://faculty.ksu.edu.sa/sites/default/files/Alfaifi_etal_2025.pdf}
\item ``Upper Faifi as an Endangered Arabic Variety: A Philological, Descriptive, and Acoustic Study'' (University of Hawai`i ScholarSpace): \url{https://scholarspace.manoa.hawaii.edu/items/59979cd6-143c-4dc0-9362-de91254522f1}
\item Alfaifi's thesis on the morphosyntax of the Faifa dialect (Ulster University): \url{https://pure.ulster.ac.uk/en/studentTheses/aspects-of-the-morphosyntax-of-the-faifa-dialect-of-saudi-arabic-/}
\item Alfaifi and Behnstedt, notes on the Jabal Fayfa dialect (Benjamins/JBE): \url{https://www.jbe-platform.com/content/books/9789027256935-sal.11.05alf}
\item Wikipedia, Faifi language: \url{https://en.wikipedia.org/wiki/Faifi_language}
\end{itemize}

\section{Religious and Community Resources}
\begin{itemize}
\item Joshua Project, Faifi in Saudi Arabia: \url{https://joshuaproject.net/people_groups/21997/SA}
\end{itemize}

\section{Digital and Audio Resources}
None found; the documentation studies above may rest on recordings, but no public archive was seen.

\section{Curriculum Development Context}
Scholarly description exists, but no pedagogical materials were found. Because schooling in Saudi Arabia uses Modern Standard Arabic and the status of Faifi relative to Arabic is unsettled, any curriculum would be supplementary, building on the recent descriptive work above.
\end{document}